\documentclass[10pt,a4paper]{amsart}
\usepackage[margin=19mm]{geometry}
\usepackage{amsmath,amssymb,mathtools}
\usepackage[scr=rsfs,cal=euler]{mathalfa}
\usepackage{xfrac}
\usepackage[unicode,hidelinks,bookmarks=false]{hyperref}
\newcommand{\ie}{i.e.\ }
\newcommand{\cf}{cf.\ }
\newcommand{\resp}{respectively\ }
\newcommand{\Subsection}{\subsection}
\providecommand{\nobreakdash}{\nobreak\hbox{-}}
\setlength{\emergencystretch}{2em}
\multlinegap0pt
\usepackage{amssymb}
\usepackage{algorithm}\usepackage{algpseudocode}
\newtheorem{assumption}{Assumption}
\newtheorem{proposition}{Proposition}
\newtheorem{lemma}{Lemma}
\newtheorem{theorem}{Theorem}
\title{Establishing an $\Omega(\sqrt{d})$ complexity lower bound for PDMP samplers and how to break it: a sub-$\sqrt{d}$ algorithm for Gaussian-tailed targets}
\author{Augustin Chevallier}
\date{Source: arXiv:2606.19909v1 (CC BY 4.0)}
\begin{document}
\maketitle

\noindent\emph{Source and licence.} Establishing an $\Omega(\sqrt{d})$ complexity lower bound for PDMP samplers and how to break it: a sub-$\sqrt{d}$ algorithm for Gaussian-tailed targets. Version arXiv:2606.19909v1 (CC BY 4.0), distributed by arXiv under the Creative Commons Attribution 4.0 International licence, http://creativecommons.org/licenses/by/4.0/, as recorded in the arXiv licence metadata for this exact version and shown on the abstract page https://arxiv.org/abs/2606.19909v1. Full licence notice: \texttt{licenses/arxiv-2606.19909-CC-BY-4.0.txt}.\par\smallskip

\noindent\textit{Editorial scope.} Counted unit: the article's theorem th:complexity (source0.tex lines 424--427). The article's complete original proof of it is delivered with it (source0.tex lines 430--474, one \texttt{proof} environment of 45 lines). No mathematics is written, repaired or re-derived: the statement and the proof are the article's own lines, in the article's own notation, and no retained line is substituted or re-flowed.  Retained uncounted as the source-local context the proof needs: lines 294--300; lines 326--332; lines 335--344; lines 357--363; lines 368--378; lines 413--422; lines 796--799.  Nothing was omitted from any retained span, so the package carries no omission record.  No retained line contains a citation, so the wrapper carries no bibliography and no bibliography entry.  No retained line contains a figure environment, an image-inclusion command or drawing code, so no image is delivered and the package declares no asset dependency.  Editorial formatting only, in addition: an AMS \texttt{amsart} preamble that restates the article's own declarations and macros used by the retained lines, namely the environments \texttt{assumption}, \texttt{proposition}, \texttt{lemma}, \texttt{theorem} on separate counters, together with the standard packages those lines need.  The retained environments print in this document as Assumption 1, Proposition 1, Lemma 1, Theorem 1 in order of appearance, whereas the article prints the same items with the numbering of its own full document.  The complete native source of the article as distributed in the arXiv e-print for this exact version is bundled unchanged as \texttt{sources/203178/source0.tex}.
\begin{assumption}[Effectively Independent Sampling Times]
    \label{assum:independent_times}
    There exists an infinite sequence of random times $0 \leq t_1 < t_2 < \dots < t_k < \dots$ increasing almost surely to infinity, such that the sequence of spatial states $(X_{t_k})_{k \geq 1}$ acts as a sequence of independent draws from $\pi$. Specifically, we assume that the empirical average of the distances between consecutive points converges almost surely to the expected distance between two independent random vectors distributed according to $\pi$:
    \[
        \lim_{K \to \infty} \frac{1}{K} \sum_{k=1}^{K-1} \|X_{t_{k+1}} - X_{t_k}\| = \mathbb{E}_{\pi \otimes \pi}[\|X - Y\|] \quad a.s.
    \]
\end{assumption}

\begin{assumption}[Ergodicity]
    \label{assum:ergodicity}
    The PDMP $(X_t, Z_t)_{t \geq 0}$ is ergodic with respect to its invariant measure $\mu$. Specifically, for any $\mu$-integrable function $f$, the time average along the trajectory converges almost surely to the spatial expectation:
    \[
        \lim_{S \to \infty} \frac{1}{S} \int_0^S f(X_t, Z_t) dt = \mathbb{E}_\mu[f] \quad a.s.
    \]
\end{assumption}

\begin{proposition}[Lower Bound on PDMP Sampling Complexity for i.i.d. targets]
    \label{prop:complexity_lower_bound}
    Let the target distribution $\pi$ on $\mathbb{R}^d$ be an i.i.d. product measure such that $\pi(x) = \prod_{i=1}^d \pi_1(x_i)$. Let the PDMP satisfy Assumption \ref{assum:ergodicity} and Assumption \ref{assum:independent_times}. 
    
    The sampling complexity $C_d$, defined as the almost sure asymptotic number of jump events required per effectively independent sample ($C_d = \lim_{K \to \infty} \frac{N(t_K)}{K}$), satisfies the following lower bound:
    \[
        C_d \geq \frac{C(\pi_1)\sqrt{d}}{\mathbb{E}_\mu[\|\phi_x\|]} \mathbb{E}_\mu[\lambda]
    \]
    where $C(\pi_1)$ is the constant defined in Lemma \ref{lem:distance_lower}.
\end{proposition}

\begin{lemma}[0 mean divergence]
    \label{lemma:zero-mean-div}
    Let $\mu$ be a probability distribution on some vector space $E$, and $x\mapsto \phi(x), x \mapsto \lambda(x), c\mapsto Q(x,\cdot)$ be a vector field, a jump rate and a jump kernel of a PDMP. If $\mu$ is invariant by the PDMP then:
    \[
        E_\mu[\frac{div(\phi \mu)}{\mu}] = 0
    \]
\end{lemma}

\begin{proposition}[lower bound of $\lambda$]
    \label{prop:lower-bound}
    Let $\mu$ be a probability distribution on some vector space $E$, and $x\mapsto \phi(x), x \mapsto \lambda(x), c\mapsto Q(x,\cdot)$ be a vector field, a jump rate and a jump kernel of a PDMP. If $\mu$ is invariant by the PDMP then for all $x$:
    \[
        \lambda(x) \geq max(0, -div(\phi \mu)(x) / \mu(x)). 
    \]
    This implies
    \[
        E_\mu[\lambda(x)] \geq \frac{1}{2} E_\mu[\frac{|div(\phi \mu)|}{\mu}].
    \]
\end{proposition}

\begin{assumption}[Independence Structure for i.i.d. Targets]
    \label{assum:independence_structure}
    First, the state space for a given dimension $d$ can be written as $(\mathbb{R} \times \mathcal{Z}_1)^d$, where $\mathcal{Z}_1$ is the auxiliary space used for targets in $\mathbb{R}$. 
    
    Second, for an i.i.d. target distribution---meaning $\pi(x) = \prod_{i=1}^d \pi_1(x_i)$---the algorithm yields a structure on each marginal space $\mathbb{R} \times \mathcal{Z}_1$ that is identical to what would be used in the 1D case. This implies:
    \begin{enumerate}
        \item The auxiliary conditional distribution completely factorizes: $\rho(x,z) = \prod_{i=1}^d \rho_1(x_i,z_i)$, leading to a joint invariant measure $\mu(x,z) = \prod_{i=1}^d \mu_1(x_i,z_i)$ with $\mu_1(x_1,z_1) = \pi_1(x_1) \rho_1(x_1,z_1)$.
        \item The deterministic vector field decomposes cleanly: $\phi(x,z) = (\phi_1(x_1,z_1), \dots, \phi_1(x_d,z_d))$, where $\phi_1$ is the exact vector field used in the 1D case.
    \end{enumerate}
\end{assumption}

\begin{lemma}[Lower Bound on Expected Distance for i.i.d. targets]
    \label{lem:distance_lower}
    Let $X, Y \in \mathbb{R}^d$ be independent random vectors drawn from a product measure $\pi = \pi_1^{\otimes d}$. Assume $\pi_1$ has finite variance and is not a Dirac mass. Then, there exists a constant $C > 0$ such that $\mathbb{E}_\pi[\|X - Y\|_2] \ge C \sqrt{d}$.
\end{lemma}

\begin{theorem}[Complexity for i.i.d. targets]
    \label{th:complexity}
    Under Assumption \ref{assum:independence_structure}, if the divergence of the 1D process has a non-zero finite variance, the expected jump rate scales as $\Omega(\sqrt{d})$, and the overall sampling complexity scales as $C = \Omega(\sqrt{d})$.
\end{theorem}

\begin{proof}
    Let $S_d(x,z) = \frac{\text{div}(\phi \mu)(x,z)}{\mu(x,z)}$. Given the factorized structure of both the vector field $\phi$ and the invariant measure $\mu$, the divergence operator distributes over the individual coordinates:
    $$ \text{div}(\phi\mu) = \sum_{i=1}^d \text{div}_1(\phi_1 \mu_1)(x_i, z_i) \prod_{j \neq i} \mu_1(x_j, z_j) $$

    Dividing by the joint measure $\mu = \prod_{j=1}^d \mu_1(x_j, z_j)$ cancels out the product terms, yielding a sum of independent ratios:
    $$ S_d(x,z) = \sum_{i=1}^d \frac{\text{div}_1(\phi_1 \mu_1)(x_i, z_i)}{\mu_1(x_i, z_i)} $$

    Let $Y_i = \frac{\text{div}_1(\phi_1 \mu_1)(x_i, z_i)}{\mu_1(x_i, z_i)}$. Under the stationary measure $\mu$, the random variables $Y_1, \dots, Y_d$ are independent and identically distributed. By Lemma \ref{lemma:zero-mean-div} applied to the 1D marginals, we know that $E_{\mu_1}[Y_i] = 0$. Let $\sigma^2 = E_{\mu_1}[Y_i^2]$ be the strictly positive, finite variance of these terms.

    By the Central Limit Theorem (CLT), as the dimension $d \to \infty$, the normalized sum $\frac{S_d}{\sigma \sqrt{d}}$ converges in distribution to a standard normal $\mathcal{N}(0, 1)$. To deduce the asymptotic behavior of the expectation $E_\mu[|S_d|]$, we must establish that the sequence $\left\{ \frac{S_d}{\sigma \sqrt{d}} \right\}_{d \ge 1}$ is uniformly integrable. 

    Because the variables $Y_i$ are independent with mean zero and variance $\sigma^2$, the variance of the normalized sum is constant for all $d$:
    $$ E_\mu\left[ \left( \frac{S_d}{\sigma \sqrt{d}} \right)^2 \right] = \frac{1}{\sigma^2 d} (d \sigma^2) = 1 $$

    Since the sequence $\left\{ \frac{S_d}{\sigma \sqrt{d}} \right\}_{d \ge 1}$ is bounded in $L^2(\mu)$, it is uniformly integrable in $L^1(\mu)$. This permits us to pass the limit inside the expectation, yielding convergence in $L^1$:
    $$ \lim_{d \to \infty} E_\mu\left[ \left| \frac{S_d}{\sigma \sqrt{d}} \right| \right] = E[|\mathcal{N}(0, 1)|] = \sqrt{\frac{2}{\pi}} $$

    Multiplying by the normalization factor $\sigma \sqrt{d}$, we obtain the exact asymptotic equivalence for the expected absolute sum:
    $$ E_\mu[|S_d|] \sim \sigma \sqrt{\frac{2}{\pi}} \sqrt{d} = \Theta(\sqrt{d}) $$

    Using the lower bound established proposition \ref{prop:lower-bound}, the expected jump rate must satisfy:
    $$ E_\mu[\lambda] \geq \frac{1}{2} E_\mu\left[ \frac{|\text{div}(\phi \mu)|}{\mu} \right] = \frac{1}{2} E_\mu[|S_d|] = \Omega(\sqrt{d}) $$

        To evaluate the overall sampling complexity $C_d$, we now invoke Proposition \ref{prop:complexity_lower_bound}, which establishes that:
    \[
        C_d \geq \frac{C \sqrt{d}}{\mathbb{E}_\mu[\|\phi_x\|]} \mathbb{E}_\mu[\lambda]
    \]
    where $C > 0$ is the distance constant from Lemma \ref{lem:distance_lower}. 
    
    We must bound the expected speed $\mathbb{E}_\mu[\|\phi_x\|]$ in the denominator. By our structural independence assumptions, the continuous flow decomposes as $\phi_x(x,z) = (\phi_{x,1}(x_1,z_1), \dots, \phi_{x,1}(x_d,z_d))$. Applying Jensen's inequality, we bound the $L^1$ norm by the $L^2$ norm:
    \[
        \mathbb{E}_\mu[\|\phi_x\|] \le \sqrt{\mathbb{E}_\mu[\|\phi_x\|^2]}  = \sqrt{d \mathbb{E}_{\mu_1}[\phi_{x,1}^2]}
    \]
    Assuming the 1D flow has a finite second moment (let $v_1^2 = \mathbb{E}_{\mu_1}[\phi_{x,1}^2] < \infty$), we have $\mathbb{E}_\mu[\|\phi_x\|] \le v_1 \sqrt{d}$.
    
    Substituting this upper bound into our complexity inequality gives:
    \[
        C_d \geq \frac{C \sqrt{d}}{v_1 \sqrt{d}} \mathbb{E}_\mu[\lambda] = \frac{C}{v_1} \mathbb{E}_\mu[\lambda]
    \]
    Since $C$ and $v_1$ are constants strictly independent of the dimension $d$, and we have already shown that $\mathbb{E}_\mu[\lambda] = \Omega(\sqrt{d})$, it immediately follows that:
    \[
        C_d = \Omega(\sqrt{d})
    \]
    This completes the proof.
\end{proof}

\end{document}
